\section{射影闭包}\label{sec:projective-closures}

上一节把射影簇看成一个仿射簇与一组无穷远点的不交并. 本节反过来, 描述怎样向仿射簇添加无穷远点得到射影簇. 令 $j_{0}:\mathbb{A}^{n}\to\mathbb{P}^{n}$ 为函数
\[
  j_{0}(a_{1},\ldots,a_{n})=[1:a_{1}:\cdots:a_{n}],
\]
它是从 $\mathbb{A}^{n}$ 到仿射开片 $\mathbb{A}_{0}^{n}\subseteq\mathbb{P}^{n}$ 的双射. 若 $X\subseteq\mathbb{A}^{n}$ 为仿射簇, 则按构造, $j_{0}(X)\subseteq\mathbb{P}^{n}$ 是 $\mathbb{P}^{n}$ 中第一个仿射限制为 $X$ 的子集, 但 $j_{0}(X)$ 一般不是射影簇. 为了把 $j_{0}(X)$ 扩充成射影簇, 必须再添加一些点. 由此得到的最小射影簇称为 $X$ 的射影闭包.

\begin{Definition}{射影闭包}\leavevmode\label{def:projective-closure}
设 $X\subseteq\mathbb{A}^{n}$ 为仿射簇, $j_{0}(X)$ 为 $X$ 在第一个仿射开片 $\mathbb{A}_{0}^{n}\subseteq\mathbb{P}^{n}$ 中的像. $X$ 的\emph{射影闭包}, 记为 $\overline{X}\subseteq\mathbb{P}^{n}$, 是所有包含 $j_{0}(X)$ 的射影簇的交.
\end{Definition}

由于射影簇的交 (即使是无穷多个射影簇的交) 仍然是射影簇 (Exercise~\ref{ex:projective-varieties-closed-under-unions-and-intersections}), $\overline{X}$ 本身就是射影簇. 按定义, $\overline{X}$ 包含于每个包含 $j_{0}(X)$ 的射影簇, 所以可以把它看成包含 $j_{0}(X)$ 的最小射影簇.

\begin{Example}{直线的射影闭包}\leavevmode\label{ex:projective-closure-of-affine-line}
若 $X=\mathcal{V}_{\mathbb{A}}(1+x_{1}-x_{2})=\{(a,1+a)\mid a\in K\}\subseteq\mathbb{A}^{2}$, 则
\[
  j_{0}(X)=\{[1:a:1+a]\mid a\in K\}\subseteq\mathbb{P}^{2}.
\]
特别地, 当 $a\neq0$ 时,
\[
  [1:a:1+a]=[\tfrac{1}{a}:1:\tfrac{1}{a}+1],
\]
所以 $j_{0}(X)$ 包含所有形如 $[b:1:b+1]$ 且 $b\in K\setminus\{0\}$ 的点. 但包含所有这些点的射影簇也必须包含 $b=0$ 时的点 (Exercise~\ref{ex:line-limit-point-belongs-to-projective-variety}). 关键是, $K[x_{0},x_{1},x_{2}]$ 中任一多项式, 如果在所有 $b\in K\setminus\{0\}$ 对应的 $(b,1,b+1)$ 上取值为零, 那么它在 $(0,1,1)$ 上也必取值为零. 由于 $[0:1:1]\notin j_{0}(X)$, $j_{0}(X)$ 不可能是射影簇.

添加缺少的点 $[0:1:1]$ 得到射影簇
\[
  \mathcal{V}_{\mathbb{P}}(x_{0}+x_{1}-x_{2})=j_{0}(X)\sqcup\{[0:1:1]\},
\]
只须按第一个坐标为零或非零, 将 $\mathcal{V}_{\mathbb{P}}(x_{0}+x_{1}-x_{2})$ 的点分成两个不交的集合即可验证. $j_{0}(X)$ 本身不是射影簇, 而 $\mathcal{V}_{\mathbb{P}}(x_{0}+x_{1}-x_{2})$ 只比它多一个点, 所以 $\mathcal{V}_{\mathbb{P}}(x_{0}+x_{1}-x_{2})$ 是包含 $j_{0}(X)$ 的最小射影簇. 因此 $\overline{X}=\mathcal{V}_{\mathbb{P}}(x_{0}+x_{1}-x_{2})$.
\end{Example}

\begin{Example}{抛物线的射影闭包}\leavevmode\label{ex:projective-closure-of-affine-parabola}
考虑抛物线
\[
  X=\mathcal{V}_{\mathbb{A}}(x_{2}-x_{1}^{2})\subseteq\mathbb{A}^{2}.
\]
那么
\[
  j_{0}(X)=\{[1:a:a^{2}]\mid a\in K\}\subseteq\mathbb{P}^{2}.
\]
和上一例一样, 当 $a\neq0$ 时,
\[
  [1:a:a^{2}]=[\tfrac{1}{a^{2}}:\tfrac{1}{a}:1],
\]
所以对任意 $b\in K\setminus\{0\}$, $j_{0}(X)$ 都包含点 $[b^{2}:b:1]$. 同样, 包含所有这些点的射影簇也必须包含点 $[0:0:1]$. 由于 $j_{0}(X)$ 不含这个点, 它不可能是射影簇.

添加缺少的点得到射影簇
\[
  \mathcal{V}_{\mathbb{P}}(x_{0}x_{2}-x_{1}^{2})=j_{0}(X)\sqcup\{[0:0:1]\},
\]
从而 $\overline{X}=\mathcal{V}_{\mathbb{P}}(x_{0}x_{2}-x_{1}^{2})$.
\end{Example}

在前两个例子中, 仿射簇 $X$ 由一个非齐次多项式 $f$ 定义, 而 $\overline{X}$ 的定义多项式可由 $f$\textquotedbl{}齐次化\textquotedbl{}得到, 即把 $f$ 的每一项乘以 $x_{0}$ 的某个幂, 得到一个齐次多项式. 再举一例, 令
\[
  f(x_{1},x_{2},x_{3})=x_{1}^{2}+x_{2}+x_{1}x_{3}^{4}.
\]
次数最高的是最后一项, 它的次数为 $5$. 把每一项乘以 $x_{0}$ 的适当幂, 使其次数为 $5$, 就对 $f$ 作了齐次化. 得到的齐次多项式是
\[
  \overline{f}(x_{0},x_{1},x_{2},x_{3})=x_{0}^{3}x_{2}^{2}+x_{0}^{4}x_{2}+x_{1}x_{3}^{4}.
\]
下列定义描述这个一般过程.

\begin{Definition}{多项式的齐次化}\leavevmode\label{def:homogenization-of-polynomial}
设 $f\in K[x_{1},\ldots,x_{n}]$ 为 $d$ 次多项式. $f$ 的\emph{齐次化}定义为
\[
  \overline{f}=x_{0}^{d}\cdot f\biggl(\frac{x_{1}}{x_{0}},\ldots,\frac{x_{n}}{x_{0}}\biggr)
  \in K[x_{0},x_{1},\ldots,x_{n}].
\]
\end{Definition}

读者应自行验证 $\overline{f}$ 是 $d$ 次齐次多项式 (Exercise~\ref{ex:homogenization-is-homogeneous-of-same-degree}), 并且这个定义和前面逐项处理的过程一致. 将 $x_{0}=1$ 代入定义, 得
\begin{equation}\label{eq:dehomogenize-by-setting-first-coordinate-one}
  \overline{f}(1,x_{1},\ldots,x_{n})=f(x_{1},\ldots,x_{n})\in K[x_{1},\ldots,x_{n}].
\end{equation}

在 Examples~\ref{ex:projective-closure-of-affine-line} 和~\ref{ex:projective-closure-of-affine-parabola} 中, 有 $X=\mathcal{V}_{\mathbb{A}}(f)$ 和 $\overline{X}=\mathcal{V}_{\mathbb{P}}(\overline{f})$, 这或许使人猜想, 若 $X=\mathcal{V}_{\mathbb{A}}(f_{1},\ldots,f_{k})$, 就有 $\overline{X}=\mathcal{V}_{\mathbb{P}}(\overline{f}_{1},\ldots,\overline{f}_{k})$. 如果成立, 那当然很方便. 遗憾的是, 从 $X$ 到 $\overline{X}$ 的过程并不总这么简单, 如下例所示.

\begin{Example}{扭三次曲线}\leavevmode\label{ex:projective-twisted-cubic}
令 $X=\mathcal{V}_{\mathbb{A}}(x_{2}-x_{1}^{2},x_{3}-x_{1}^{3})$, 那么
\[
  j_{0}(X)=\{[1:a:a^{2}:a^{3}]\mid a\in K\}\subseteq\mathbb{P}^{3}.
\]
与 Examples~\ref{ex:projective-closure-of-affine-line} 和~\ref{ex:projective-closure-of-affine-parabola} 类似, 任何包含 $j_{0}(X)$ 的射影簇也必须包含点 $[0:0:0:1]$. 因此 $j_{0}(X)$ 不是射影簇. 但直接计算 (Exercise~\ref{ex:equations-for-projective-twisted-cubic}) 表明 $j_{0}(X)\sqcup\{[0:0:0:1]\}$ 是射影簇:
\[
  \overline{X}=\mathcal{V}_{\mathbb{P}}(x_{0}x_{2}-x_{1}^{2},x_{0}^{2}x_{3}-x_{1}^{3},x_{1}x_{3}-x_{2}^{2})
  =j_{0}(X)\cup\{[0:0:0:1]\}.
\]
这个射影簇称为\emph{扭三次曲线}.

$\overline{X}$ 的前两个定义多项式由 $X$ 的定义多项式齐次化得到, 但第三个也是必不可少的. 若没有它, 得到的是
\begin{equation}\label{eq:extra-projective-line-in-two-equation-cubic}
  \mathcal{V}_{\mathbb{P}}(x_{0}x_{2}-x_{1}^{2},x_{0}^{2}x_{3}-x_{1}^{3}),
\end{equation}
它包含 $j_{0}(X)$, 但还包含多余的点 $[0:0:b:c]\in\mathbb{P}^{3}$. 因此, 尽管 \eqref{eq:extra-projective-line-in-two-equation-cubic} 是包含 $j_{0}(X)$ 的射影簇, 它却不是 $\overline{X}$.
\end{Example}

Example~\ref{ex:projective-twisted-cubic} 说明, $X=\mathcal{V}_{\mathbb{A}}(f_{1},\ldots,f_{r})$ 的射影闭包一般不能只对 $f_{1},\ldots,f_{r}$ 齐次化得到. 但有一个解决办法: 不要只齐次化任意选取的一组定义多项式, 而应齐次化 $X$ 的消没理想中的每个多项式. 下列记号将很有用.

\begin{Definition}{集合的齐次化}\leavevmode\label{def:homogenization-of-set}
设 $S\subseteq K[x_{1},\ldots,x_{n}]$ 为多项式集合. $S$ 的\emph{齐次化}是集合
\[
  \overline{S}=\{\overline{f}\mid f\in S\}\subseteq K[x_{0},\ldots,x_{n}].
\]
\end{Definition}

\BookPairedNote
  {将仿射簇的消没理想中的每个多项式齐次化后, 就得到足够多的多项式来描述它的射影闭包. 下一个结果就是这一陈述.}
  {注意: 即使 $I$ 是理想, 它的齐次化 $\overline{I}$ 一般也不是理想. 你能看出原因吗?}

\begin{Proposition}{通过齐次化求射影闭包}\leavevmode\label{prop:projective-closure-by-homogenizing-vanishing-ideal}
若 $X\subseteq\mathbb{A}^{n}$ 为仿射簇, 则
\[
  \overline{X}=\mathcal{V}_{\mathbb{P}}\bigl(\overline{\mathcal{I}_{\mathbb{A}}(X)}\bigr)\subseteq\mathbb{P}^{n}.
\]
\end{Proposition}
\begin{Proof}
我们证明两个包含关系.

($\subseteq$) 若 $[1:a_{1}:\cdots:a_{n}]\in j_{0}(X)$ 且 $g\in\overline{\mathcal{I}_{\mathbb{A}}(X)}$, 则对某个 $f\in\mathcal{I}_{\mathbb{A}}(X)$, 有 $g=\overline{f}$, 因而 $g(1,a_{1},\ldots,a_{n})=f(a_{1},\ldots,a_{n})=0$. 所以 $\mathcal{V}_{\mathbb{P}}\bigl(\overline{\mathcal{I}_{\mathbb{A}}(X)}\bigr)$ 是包含 $j_{0}(X)$ 的射影簇. 由于 $\overline{X}$ 是包含 $j_{0}(X)$ 的最小射影簇, 得 $\overline{X}\subseteq\mathcal{V}_{\mathbb{P}}\bigl(\overline{\mathcal{I}_{\mathbb{A}}(X)}\bigr)$.

($\supseteq$) 先证明
\begin{equation}\label{eq:projective-ideal-inclusion-for-homogenized-affine-ideal}
  \mathcal{I}_{\mathbb{P}}(\overline{X})\subseteq
  \mathcal{I}_{\mathbb{P}}\Bigl(\mathcal{V}_{\mathbb{P}}\bigl(\overline{\mathcal{I}_{\mathbb{A}}(X)}\bigr)\Bigr).
\end{equation}
对这个包含关系的两边应用 $\mathcal{V}_{\mathbb{P}}$, 由 Proposition~\ref{prop:basic-properties-of-projective-vanishing-operators} 得
\[
  \overline{X}\supseteq\mathcal{V}_{\mathbb{P}}\bigl(\overline{\mathcal{I}_{\mathbb{A}}(X)}\bigr).
\]
为证明 \eqref{eq:projective-ideal-inclusion-for-homogenized-affine-ideal}, 令 $g\in\mathcal{I}_{\mathbb{P}}(\overline{X})$. 由于 $\mathcal{I}_{\mathbb{P}}(\overline{X})$ 是齐次理想, 因而有齐次生成元, 所以为证明 \eqref{eq:projective-ideal-inclusion-for-homogenized-affine-ideal}, 只须假设 $g$ 是齐次的.

即使 $g$ 是齐次的, 它也未必是 $K[x_{1},\ldots,x_{n}]$ 中某个元素的齐次化, 因为 $g$ 的每一项都可能含有 $x_{0}$. 不过, 若 $k\in\mathbb{N}$ 是使得对某个 $h\in K[x_{0},x_{1},\ldots,x_{n}]$ 有 $g=x_{0}^{k}h$ 的最大幂指数, 那么 $h=\overline{h_{0}}$, 其中 $h_{0}=h(1,x_{1},\ldots,x_{n})\in K[x_{1},\ldots,x_{n}]$ (Exercise~\ref{ex:recover-homogeneous-polynomial-not-divisible-by-first-coordinate}). 由于 $g$ 在 $\overline{X}\supseteq j_{0}(X)$ 上取值为零, 所以 $h_{0}$ 在 $X$ 上取值为零: 对任意 $(a_{1},\ldots,a_{n})\in X$, 有
\[
  h_{0}(a_{1},\ldots,a_{n})=h(1,a_{1},\ldots,a_{n})=g(1,a_{1},\ldots,a_{n})=0.
\]
因此 $h_{0}\in\mathcal{I}_{\mathbb{A}}(X)$, 所以 $h\in\overline{\mathcal{I}_{\mathbb{A}}(X)}$, 进而
\[
  g\in\langle\overline{\mathcal{I}_{\mathbb{A}}(X)}\rangle
    \subseteq\mathcal{I}_{\mathbb{P}}\Bigl(\mathcal{V}_{\mathbb{P}}\bigl(\overline{\mathcal{I}_{\mathbb{A}}(X)}\bigr)\Bigr),
\]
其中包含关系是 Proposition~\ref{prop:basic-properties-of-projective-vanishing-operators} 的再次应用. 这就证明了 \eqref{eq:projective-ideal-inclusion-for-homogenized-affine-ideal}, 也就证明了这个命题.
\end{Proof}

\begin{Example}{再看扭三次曲线}\leavevmode\label{ex:homogenizing-twisted-cubic-ideal}
根据 Proposition~\ref{prop:projective-closure-by-homogenizing-vanishing-ideal}, 可以进一步理解 Example~\ref{ex:projective-twisted-cubic} 中的现象. 当时为了计算
\[
  X=\mathcal{V}_{\mathbb{A}}(x_{2}-x_{1}^{2},x_{3}-x_{1}^{3})
\]
的射影闭包, 我们发现只齐次化 $X$ 的两个定义多项式是不够的. Proposition~\ref{prop:projective-closure-by-homogenizing-vanishing-ideal} 保证, 齐次化 $\mathcal{I}(X)$ 中的所有多项式后, 就能得到一组定义 $\overline{X}$ 的多项式. 在这个特定例子中, 只要添上另一个多项式 $x_{1}x_{3}-x_{2}^{2}\in\mathcal{I}(X)$ 的齐次化就足够了 (这个多项式碰巧是齐次的).
\end{Example}

由扭三次曲线的例子可见, 一般并没有
\[
  \overline{\mathcal{V}_{\mathbb{A}}(f_{1},\ldots,f_{k})}
    =\mathcal{V}_{\mathbb{P}}(\overline{f}_{1},\ldots,\overline{f}_{k}).
\]
但在 $k=1$ 的特殊情形, 零点定理蕴含, 射影闭包确实可由唯一的定义多项式齐次化得到.

\begin{Proposition}{超曲面的射影闭包}\leavevmode\label{prop:projective-closure-of-hypersurface}
若 $f\in K[x_{1},\ldots,x_{n}]$, 则 $\overline{\mathcal{V}_{\mathbb{A}}(f)}=\mathcal{V}_{\mathbb{P}}(\overline{f})$.
\end{Proposition}
\begin{Proof}
令 $X=\mathcal{V}_{\mathbb{A}}(f)$, 并令 $f=q_{1}^{k_{1}}\cdots q_{\ell}^{k_{\ell}}$ 为 $f$ 的由互异不可约因子给出的分解. 零点定理蕴含 $\mathcal{I}_{\mathbb{A}}(X)=\langle q_{1}\cdots q_{\ell}\rangle$. 因此由 Proposition~\ref{prop:projective-closure-by-homogenizing-vanishing-ideal}, 有
\[
  \overline{\mathcal{V}_{\mathbb{A}}(f)}=
  \mathcal{V}_{\mathbb{P}}\bigl(\overline{\mathcal{I}_{\mathbb{A}}(X)}\bigr)=
  \mathcal{V}_{\mathbb{P}}\bigl(\overline{\langle q_{1}\cdots q_{\ell}\rangle}\bigr).
\]
还须证明
\[
  \mathcal{V}_{\mathbb{P}}\bigl(\overline{\langle q_{1}\cdots q_{\ell}\rangle}\bigr)
  =\mathcal{V}_{\mathbb{P}}(\overline{f}).
\]
由于 $f\in\langle q_{1}\cdots q_{\ell}\rangle$, 有 $\overline{f}\in\overline{\langle q_{1}\cdots q_{\ell}\rangle}$, 所以 $\mathcal{V}_{\mathbb{P}}\bigl(\overline{\langle q_{1}\cdots q_{\ell}\rangle}\bigr)\subseteq\mathcal{V}_{\mathbb{P}}(\overline{f})$. 为证明反向包含, 设 $a\in\mathcal{V}_{\mathbb{P}}(\overline{f})$, 即 $\overline{f}(a)=0$. 利用齐次化的乘法性 (Exercise~\ref{ex:homogenization-multiplication-and-addition}), 有
\[
  \overline{f}=\overline{q}_{1}^{k_{1}}\cdots\overline{q}_{\ell}^{k_{\ell}},
\]
所以 $\overline{f}(a)=0$ 蕴含对某个 $i$, 有 $\overline{q}_{i}(a)=0$. 现在任取 $g\in\overline{\langle q_{1}\cdots q_{\ell}\rangle}$, 存在 $h\in K[x_{1},\ldots,x_{n}]$, 使得
\[
  g=\overline{h\cdot q_{1}\cdots q_{\ell}}=\overline{h}\cdot\overline{q}_{1}\cdots\overline{q}_{\ell}.
\]
由于某个 $\overline{q}_{i}(a)=0$, 得 $g(a)=0$. 又因为 $g$ 是 $\overline{\langle q_{1}\cdots q_{\ell}\rangle}$ 的任一元素, 所以 $a\in\mathcal{V}_{\mathbb{P}}\bigl(\overline{\langle q_{1}\cdots q_{\ell}\rangle}\bigr)$. 因此 $\mathcal{V}_{\mathbb{P}}(\overline{f})\subseteq\mathcal{V}_{\mathbb{P}}\bigl(\overline{\langle q_{1}\cdots q_{\ell}\rangle}\bigr)$, 证明完成.
\end{Proof}

关于射影闭包的最后一个重要结果是, 取射影闭包绝不会在仿射空间内部添加点, 只会添加无穷远点. 换句话说, 取仿射限制在某种意义上是取射影闭包的逆过程. 下一个结果精确地说明这一点.

\begin{Proposition}{$(\overline{X})_{0}=X$}\leavevmode\label{prop:affine-restriction-of-projective-closure}
若 $X\subseteq\mathbb{A}^{n}$ 为仿射簇, 则 $X$ 的射影闭包的仿射限制就是 $X$. 简写为 $(\overline{X})_{0}=X$.
\end{Proposition}
\begin{Proof}
设 $X\subseteq\mathbb{A}^{n}$ 为仿射簇, 令 $S=\overline{\mathcal{I}_{\mathbb{A}}(X)}$. 结合 Proposition~\ref{prop:projective-closure-by-homogenizing-vanishing-ideal} 与 Definition~\ref{def:affine-patches-and-restrictions}, 可知射影闭包 $\overline{X}=\mathcal{V}_{\mathbb{P}}(S)$ 的仿射限制是 $\mathcal{V}_{\mathbb{A}}(S_{0})$, 其中
\[
  \begin{aligned}
    S_{0}&=\{g(1,x_{1},\ldots,x_{n})\mid g\in\overline{\mathcal{I}_{\mathbb{A}}(X)}\}\\
         &=\{\overline{f}(1,x_{1},\ldots,x_{n})\mid f\in\mathcal{I}_{\mathbb{A}}(X)\}\\
         &=\mathcal{I}_{\mathbb{A}}(X).
  \end{aligned}
\]
因此 $\overline{X}$ 的仿射限制是 $\mathcal{V}_{\mathbb{A}}(\mathcal{I}_{\mathbb{A}}(X))=X$, 如所断言.
\end{Proof}

仿射限制和射影闭包在仿射簇与射影簇之间建立了密切的联系. 如前所见, 两种情形下的理论基本上是平行的. 因此, 读者或许会预期零点定理也有射影版本. 事实确实如此, 本章最后一节将证明它.

\BookExerciseHeading
\begin{BookExercises}
  \item\label{ex:line-limit-point-belongs-to-projective-variety} 设 $X\subseteq\mathbb{P}^{2}$ 为射影簇, 且对所有 $b\in K\setminus\{0\}$ 都包含点 $[b:1:b+1]$. 证明 $[0:1:1]\in X$.
  \item\label{ex:projective-variety-contains-parameter-zero-limit} 推广上一题. 设 $f_{0},\ldots,f_{n}\in K[x]$ 为一元多项式, 它们在 $0$ 处的取值不全为零. 设 $X\subseteq\mathbb{P}^{n}$ 为射影簇, 使得对所有 $b\in K\setminus\{0\}$, 都有
  \[
    [f_{0}(b):\cdots:f_{n}(b)]\in X.
  \]
  证明
  \[
    [f_{0}(0):\cdots:f_{n}(0)]\in X.
  \]
  \item\label{ex:projective-closure-of-affine-hyperbola} 改用 Examples~\ref{ex:projective-closure-of-affine-line} 和~\ref{ex:projective-closure-of-affine-parabola} 的论证, 证明仿射簇
  \[
    X=\mathcal{V}_{\mathbb{A}}(x_{1}x_{2}-1)\subseteq\mathbb{A}^{2}
  \]
  的射影闭包必须包含点 $[0:0:1]$ 和 $[0:1:0]$, 并推出
  \[
    \overline{X}=\mathcal{V}_{\mathbb{P}}(x_{1}x_{2}-x_{0}^{2}).
  \]
  画出 $X$ 在 $\mathbb{R}$ 上的图形. $X$ 的渐近线是什么, 它们与 $\overline{X}$ 的无穷远点有什么关系?
  \item\label{ex:equations-for-projective-twisted-cubic} 令
  \[
    X=\mathcal{V}_{\mathbb{A}}(x_{2}-x_{1}^{2},x_{3}-x_{1}^{3})\subseteq\mathbb{A}^{3}
  \]
  为仿射扭三次曲线.
  \begin{enumerate}
    \item 证明 $\overline{X}=\mathcal{V}_{\mathbb{P}}(x_{0}x_{2}-x_{1}^{2},x_{0}^{2}x_{3}-x_{1}^{3},x_{1}x_{3}-x_{2}^{2})$.
    \item 证明 $\overline{X}=\mathcal{V}_{\mathbb{P}}(x_{0}x_{2}-x_{1}^{2},x_{0}x_{3}-x_{1}x_{2},x_{1}x_{3}-x_{2}^{2})$.
  \end{enumerate}
  ((b) 中给出的描述是扭三次曲线更常见的表示. 注意, (b) 的三个方程是行向量为 $(x_{0},x_{1},x_{2})$ 和 $(x_{1},x_{2},x_{3})$ 的 $2\times3$ 矩阵的 $2\times2$ 子式.)
  \item\label{ex:homogenization-is-homogeneous-of-same-degree} 证明, 对任意 $d$ 次多项式 $f\in[x_{1},\ldots,x_{n}]$, 齐次化 $\overline{f}$ 是 $d$ 次齐次多项式.
  \item\label{ex:homogenization-multiplication-and-addition} 令 $f,g\in K[x_{1},\ldots,x_{n}]$.
  \begin{enumerate}
    \item 证明 $\overline{f\cdot g}=\overline{f}\cdot\overline{g}$.
    \item 是否有 $\overline{f+g}=\overline{f}+\overline{g}$? 给出证明或反例.
  \end{enumerate}
  \item\label{ex:recover-homogeneous-polynomial-not-divisible-by-first-coordinate} 设 $h\in K[x_{0},\ldots,x_{n}]$ 为齐次多项式, 且 $x_{0}\nmid h$. 证明 $h$ 是
  \[
    h(1,x_{1},\ldots,x_{n})\in K[x_{1},\ldots,x_{n}]
  \]
  的齐次化.
  \item\label{ex:closure-of-affine-ideal-variety-via-homogenization} 对任意理想 $I\subseteq K[x_{1},\ldots,x_{n}]$, 证明 $\overline{\mathcal{V}_{\mathbb{A}}(I)}=\mathcal{V}_{\mathbb{P}}(\overline{I})$. 提示: 用零点定理和 Proposition~\ref{prop:projective-closure-by-homogenizing-vanishing-ideal}, 只须证明
  \[
    \mathcal{V}_{\mathbb{P}}\bigl(\overline{\sqrt{I}}\bigr)=\mathcal{V}_{\mathbb{P}}(\overline{I}).
  \]
\end{BookExercises}
